% CRON_SECTION_h2logh.tex --- merge-ready block for the %%MERGE-SLOT in proof.tex
% Theorem [ABSTRACT-Q-H2LOGH]; triple-verified (ledger AT.9: hand-check + machine
% lower bound + independent hostile audit Q6620 CONFIRMED). Source: Q6592 + Q6637.
% Single-writer: cron. Life: pull this content into the merge slot verbatim.

\begin{theorem}[Abstract inverse-square clique bound]\label{thm:abstract-h2logh}
Let $x_0,\ldots,x_N$ be a word over an arbitrary alphabet. Put
$R_d=\#\{r:x_r=x_{r+d}\}$ and
$d_H(r)=\#\{1\leq h\leq H:x_r=x_{r+h}\}$. If $R_d\leq Cd$ for
$1\leq d\leq H$, then
$Q_H=\sum_r\binom{d_H(r)}2\leq22CH^2(1+\log H)$.
\end{theorem}

\begin{lemma}[Clique energy]\label{lem:clique-energy}
If $A$ is a color class inside an interval of at most $2H$ consecutive
integers, $m=|A|$, and
$E_H(A)=\sum_{u<v\in A,v-u\leq H}(v-u)^{-2}$, then
$\binom m3\leq11H^2E_H(A)$.
\end{lemma}

\begin{proof}
For $m<3$ this is immediate. If $m=3$, write the two consecutive gaps as
$g_1,g_2$. Their sum is at most $2H-1$, hence one is at most $H-1$.
Therefore $E_H(A)\geq H^{-2}$ and the claim follows.

Assume $m\geq4$. Split the containing interval into two halves. One half
contains $n\geq\lceil m/2\rceil$ points and has span at most $H-1$.
Let its consecutive gaps be $g_1,\ldots,g_k$, $k=n-1$. Cauchy gives
\[
\sum_i g_i^{-2}\geq \frac{(\sum_i g_i^{-1})^2}{k},\qquad
(\sum_i g_i)(\sum_i g_i^{-1})\geq k^2.
\]
Since $\sum_i g_i\leq H-1<H$,
$\sum_i g_i^{-2}\geq k^3/H^2$. The boundary cases $m=4,5$ give
$k\geq m/4$, and the same inequality follows for all $m\geq4$.
Thus $E_H(A)\geq m^3/(64H^2)$, and
$\binom m3\leq m^3/6\leq(64/6)H^2E_H(A)<11H^2E_H(A)$.
\end{proof}

\begin{lemma}[Window localization]\label{lem:window-localize}
For blocks $B_j=[jH,(j+1)H)$ and windows $W_j=[jH,(j+2)H)$, the triple
count satisfies
\[
\sum_{j,a}E_H(A_{j,a})\leq2\sum_{d\leq H}R_d/d^2\leq2C(1+\log H).
\]
\end{lemma}

\begin{proof}
A term of $Q_H$ is a monochromatic triple $u<v<w$ with $w-u\leq H$.
It is counted exactly once by the least point $u$ and the pair of gaps
$(v-u,w-u)$. If $u\in B_j$, then $w\in W_j$, so the whole triple lies in
that window. A pair at distance $d\leq H$ with left endpoint in $B_t$ can
occur only in $W_{t-1}$ or $W_t$, hence has multiplicity at most two.
Summing pair energies and using the harmonic series gives the claim.
\end{proof}

\begin{proof}[Proof of the theorem]
Assign each monochromatic triple (equivalently, each term of $Q_H$) to the
block of its least point; the triple then lies in the corresponding window, so
\[
Q_H\leq\sum_{j,a}\binom{m_{j,a}}3\leq11H^2\sum_{j,a}E_H(A_{j,a}),
\]
by Lemma~\ref{lem:clique-energy} applied to every window $W_j$ (an interval of
at most $2H$ consecutive positions) and color class $A_{j,a}$. The
localization lemma bounds the total energy by $2C(1+\log H)$. Hence
\[
Q_H\leq 11H^2\cdot2C(1+\log H)=22CH^2(1+\log H).
\]
\end{proof}

\begin{proposition}[Sharpness]\label{prop:word-lower}
There are words with $R_d\leq(5/2)d$ and
$Q_H=(1/8+o(1))H^2$.
\end{proposition}

\begin{proof}
Let $M=\lfloor H/2\rfloor$. For each $1\leq s\leq t\leq M$ put a private
symbol at $b_e,b_e+t,b_e+2t$, where $b_e=e(H+1)$, and fill all other
positions with fresh singleton symbols. Then
\[
R_d=2n_d+\mathbf1_{2\mid d}n_{d/2},\qquad n_t=t,
\]
so $R_d\leq5d/2$. Each class gives one triple of span $2t\leq H$, hence
$Q_H=M(M+1)/2=(1/8+o(1))H^2$. Cross-block near coincidences are not
collisions because symbols are private.
\end{proof}

\begin{corollary}\label{cor:apery-h2}
For the Ap\'ery projective orbit one has $R_d\leq 3(d-1)$; with $C=3$ and
$H=\min(\Delta,p-3)$,
$Q_p(\Delta)\leq66H^2(1+\log H)$.
\end{corollary}

\begin{corollary}\label{cor:t3-log}
$T_3=6Q_p(p-2)\leq396(p-2)^2(1+\log(p-2))$.
\end{corollary}

\begin{remark}
The exponent $2$ is optimal in the abstract class. Removing the logarithm
would follow from a prefix estimate $S_1(Y)=\sum_{d\leq Y}R_d\ll
Y^{2-\delta}$: Abel summation gives
$\sum_{d\leq H}R_d/d^2\ll1+\int_1^H S_1(t)t^{-3}dt$, which is bounded when
$S_1(t)\ll t^{2-\delta}$.
\end{remark}

